% БҮЛЭГ 8 — ДҮГНЭЛТ

\chapter{Дүгнэлт}
\label{Chapter8}

%-------------------------------------------------------------------------------
\section{Бүлэг бүрийн оруулсан хувь нэмэр}

Энэхүү судалгааны ажил нь Монгол хэл дээрх цахим орчны хортой сөрөг болон бүтээлч шүүмжлэлийг ялгах асуудлыг томьёолж (бүлэг~\ref{Chapter1}--\ref{Chapter2}), MN-BERT загварыг нарийн тохируулах (fine-tuning) арга зүйг боловсруулсан (бүлэг~\ref{Chapter3},~\ref{Chapter5}). Нийт $10{,}000$ дээжийг $\kappa{=}0.72$ нийцэлтэйгээр 4 ангилалд шошголсон (бүлэг~\ref{Chapter4}); эцсийн нэг шатлалт загвар $\mathbf{83.26\%}$ нарийвчлал, $\mathbf{0.8062}$ макро-F1 үзүүлэлт гаргаж өмнөх судалгааны нарийвчлалын үзүүлэлтээс ($82.70\%$) бага хэмжээгээр давсан (бүлэг~\ref{Chapter6}).

%-------------------------------------------------------------------------------
\section{Зааварчилгааны шаардлагын биелэлт}

Заавал биелүүлэх Ш-1-ээс Ш-7 шаардлагыг үнэлгээний рубрикийн (guideline §10)
жишиг онооны хамт, мөн судалгааны чиглэлд тохирох сонгомол НШ шаардлагуудыг
дараах хүснэгтэд нэгтгэв.

\begin{table}[H]
\centering
\caption{Шаардлагын биелэлтийн нэгдсэн хүснэгт}
\label{tab:requirements}
{\small\setlength{\tabcolsep}{4pt}
\begin{tabular}{@{}p{0.9cm}p{3.1cm}p{5.5cm}p{2.7cm}@{}}
\toprule
\textbf{Код} & \textbf{Шаардлага} & \textbf{Биелэлтийн нотолгоо} & \textbf{Төлөв} \\
\midrule
Ш-1 & Загварын гүйцэтгэл & Бүлэг~\ref{Chapter6}: нэг шатлалт MN-BERT нарийвчлал $83.26\%$ ($80$--$89\%$ муж) & 3 оноо \\
\addlinespace[2pt]
Ш-2 & Өгөгдлийн чанарын сайжруулалт & Олон үе шаттай дахин шошгололт (Хүснэгт~\ref{tab:dq-progression}): нарийвчлал $0.7253{\to}0.8326$ ($\geq$10\%) & 4 оноо \\
\addlinespace[2pt]
Ш-3 & Архитектурын шинэлэг байдал & Нэг шатлалт MN-BERT, хоёр шатлалт нэгдсэн дамжлага болон хуваалцсан энкодертой хувилбарын харьцуулсан судалгаа, MN-BERT нарийн тохируулга (Бүлэг~\ref{Chapter3},~\ref{Chapter5}) & 4 оноо \\
\addlinespace[2pt]
Ш-4 & Оновчлол ($\geq$20\% бууралт) & Нэг шатлалт загвар AMD DirectML (GPU) дээр $8.7$~мс дундаж хурд, хоёр шатлалтаас $30.8\%$ хурдан; бодит цагийн ${\leq}100$~мс шаардлага биелсэн (Хүснэгт~\ref{tab:latency}) & Биелэв \\
\addlinespace[2pt]
Ш-5 & Ёс зүй, bias, privacy шинжилгээ & Бүлэг~\ref{Chapter7} бүхэлдээ & Бүрэн \\
\addlinespace[2pt]
Ш-6 & Дахин давтагдах туршилтын орчин & random seed, орчин, код, лог (Бүлэг~\ref{Chapter5}) & Давтагдахуйц \\
\addlinespace[2pt]
Ш-7 & Академик текст ба хамгаалалт & Энэ бүхий л дипломын текст, хамгаалалт & Сайн \\
\midrule
\multicolumn{4}{@{}p{12.2cm}@{}}{\textit{Сонгомол нэмэлт шаардлага (судалгааны чиглэлд тохирохыг сонгов):}}\\
\addlinespace[2pt]
НШ-1 & Суурь (baseline) загвар байгуулах & LR/NB/SVM суурь загвар, тоон утга тогтоосон (Бүлэг~\ref{Chapter6}) & $\checkmark$ \\
\addlinespace[2pt]
НШ-2 & $\geq\!3$ алгоритмын харьцуулалт & Хүснэгт~\ref{tab:comparison-summary} (суурь, хоёр шатлалт, хуваалцсан энкодертой, нэг шатлалт) & $\checkmark$ \\
\addlinespace[2pt]
НШ-11 & Абляцийн судалгаа & Хоёр шатлалт шийдвэрлэлтийн дүрмийн туршилт (Хүснэгт~\ref{tab:two-stage-rules}), чанаржуулалтын явц (Хүснэгт~\ref{tab:dq-progression}) & $\checkmark$ \\
\addlinespace[2pt]
НШ-14 & Нийгмийн нөлөөллийн үнэлгээ & Нийгмийн нөлөөллийн үнэлгээ & $\checkmark$ \\
\bottomrule
\end{tabular}}
\end{table}

\textit{Тэмдэглэл:} Ш-1 шаардлагын хувьд $80$--$89\%$ нь 3 оноо авах бөгөөд эцсийн загвар нь суурь загваруудын хамгийн сайн үзүүлэлтээс ($66.91\%$) $+16.35$ нэгжээр илүү байна. Ш-4 шаардлага AMD DirectML (GPU) дээр нэг шатлалт загварын дундаж $8.7$~мс хурдтайгаар (ХО-19: ${\leq}100$~мс биелэв; ХО-5: нэг шатлалт нь хоёр шатлалтаас $30.8\%$ хурдан) биелэгдэв.

%-------------------------------------------------------------------------------
\section{Хязгаарлалт}

Судалгааны хязгаарлалтыг дараах байдлаар тодорхойлж байна:

\begin{enumerate}
  \item \textbf{Өгөгдлийн цар хүрээ:} 10,000 дээж бүхий өгөгдлийн багц нь
    тусгай нэр томьёоны (эрүүл мэнд, санхүү, хууль) хамралтыг бүрэн
    илэрхийлж чадахгүй байж болно.
  \item \textbf{Платформын хамаарал:} Гурван эх сурвалжийн ихэнх хэсэг нь ($\approx$91\%)
    мэдээний сайтууд бөгөөд TikTok, YouTube, Twitter/X зэрэг платформ ороогүй
    нь загвар шинэ платформын текст дээр тогтвортой ажиллахад сөрөг нөлөөлж болзошгүй.
  \item \textbf{Егөөдөл, контекстийн алдаа:} MN-BERT загвар нь егөөдсөн илэрхийлэл болон давхар утгыг танихдаа өгүүлбэрийн өнгөн талын шинж тэмдгүүдэд тулгуурлаж байгаа нь ажиглагдсан (бүлэг~\ref{Chapter6}).
    \item \textbf{Туршилтын үр дүн:} Зөвхөн нэг санамсаргүй утга (random seed $42$) болон нэг удаагийн тусгаарлах (hold-out) аргыг ашигласан тул цаашид k-fold давтан туршилт хийх нь зүйтэй.
  \end{enumerate}

  %-------------------------------------------------------------------------------
  \section{Цаашдын судалгааны чиглэл}

  \begin{enumerate}
    \item \textbf{Өгөгдлийг өргөжүүлэх:} TikTok, YouTube, Twitter/X болон форумуудаас нэмэлт өгөгдөл цуглуулж, $\geq$50,000 шошготой Монгол хэлний хортой агуулгыг илрүүлэх өгөгдлийн багцыг нийтэд нээлттэй болгох.
    \item \textbf{Том хэлний загварууд:} mBERT, XLM-R болон GPT-д суурилсан few-shot/instruction нарийн тохируулга (fine-tuning) хийж, гүйцэтгэлийн дээд хязгаарыг тодорхойлох.
    \item \textbf{Тайлбарлах боломжтой байдал (explainability):} SHAP эсвэл анхаарлын механизм (attention mechanism) ашиглан таамаглалын үндэслэлийг тайлбарлах; егөөдлийг илрүүлэх (sarcasm detection) асуудлыг олон даалгаварт сургалт (multi-task learning) аргаар нэгтгэх.
    \item \textbf{Хоёр шатлалт хувилбарын хамтарсан сургалт (end-to-end training):} Хоёр шатлалт дамжлагын хэсгүүдийг хамтад нь сургах. Одоогийн нэг шатлалт загвар GPU дээр $8.7$~мс дундаж хурдтай бөгөөд бодит цагийн шаардлагыг хангаж байна; CPU горимд INT8 квантчилал (quantization) болон дистилляци (distillation) аргуудыг ашиглан хурдыг цаашид нэмэгдүүлэх боломжтой.
    \item \textbf{API ба аудитын хэрэглээ:} FastAPI прототипийг бодит сэтгэгдлийн урсгалтай холбож, \texttt{/predict}, \texttt{/batch}, \texttt{/audit} төгсгөлүүдээр хүний хяналттай аудитын ажлын урсгалд турших.
  \end{enumerate}

  \section{Нийгмийн нөлөөллийн үнэлгээ}

  Монгол хэл дээрх цахим орчинд хортой агуулга нэмэгдэх нь хэрэглэгчийн сэтгэл зүй, онлайн орчны аюулгүй байдал, хэлэлцүүлгийн чанарт сөрөг нөлөө үзүүлдэг. Энэхүү систем нь хортой дайралтыг илрүүлэхийн зэрэгцээ үйлчилгээний чанар, байгууллагын үйл ажиллагаа, системийн доголдол болон хийсэн ажил, шийдвэрт чиглэсэн бүтээлч шүүмжлэлийг тусад нь таних боломжийг дэмжинэ. Гэхдээ систем нь зөвхөн ``зөвлөх'' үүрэгтэй байх бөгөөд эцсийн шийдвэрийг хүн гаргах ёстой гэж үзэж байна.

  Энэхүү ажил нь Монгол хэлний боловсруулалтад (NLP) хортой сөрөг болон бүтээлч шүүмжлэлийг ялгах системчилсэн нэгэн шийдэл, нийтэд нээлттэй 4 ангилалт өгөгдлийн багц болон $\mathbf{83.26\%}$ нарийвчлал ($\mathbf{0.8062}$ макро-F1) гэсэн чухал үр дүнгүүдийг гаргаж байна.
